% !PRINTED 251
\weberpage{251}
%<<CONT
jenen Ausdrücken $\alpha$ mit $\beta$ vertauschen, so gehen $\psi_1$, $\psi_2$, $\psi_3$ in $\psi_1$, $-\psi_3$, $-\psi_2$ über, und wenn wir gleichzeitig $\alpha$ mit $\beta$ und $\gamma$ mit $\delta$ vertauschen, $\psi_1$, $\psi_2$, $\psi_3$ in $\psi_1$, $-\psi_2$, $-\psi_3$. Daraus ergiebt sich Folgendes:

Setzt man für die Variable $x$ einen beliebigen reellen Werth und sind die Wurzeln $\alpha$, $\beta$, $\gamma$, $\delta$ der biquadratischen Form alle reell, so sind auch $\psi_1$, $\psi_2$, $\psi_3$ reell.

Sind $\alpha$, $\beta$ conjugirt imaginär, $\gamma$, $\delta$ reell, so ist die Vertauschung von $\alpha$ und $\beta$ gleichbedeutend mit der Vertauschung von $i$ mit $-i$. Es sind also in diesem Falle $\psi_1$ reell, $\psi_2$, $\psi_3$ conjugirt imaginär.

Sind endlich $\alpha$, $\beta$ und $\gamma$, $\delta$ zwei Paare conjugirt imaginärer Wurzeln, so werden $\alpha$ mit $\beta$ und $\gamma$ mit $\delta$ vertauscht, wenn $i$ in $-i$ übergeht; also sind in diesem Falle $\psi_1$ reell, $\psi_2$, $\psi_3$ \gesperrt{rein imaginär}.

Sind vier Wurzeln reell, so werden $\psi_1^2$, $\psi_2^2$, $\psi_3^2$ reell und positiv.

Sind zwei Wurzeln reell, so ist $\psi_1^2$ reell und positiv, $\psi_2^2$, $\psi_3^2$ sind conjugirt imaginär.

Sind vier Wurzeln imaginär, so ist $\psi_1^2$ positiv, $\psi_2^2$, $\psi_3^2$ sind reell und negativ.

Nun haben wir im \S.\,66, (7) die cubische Gleichung aufgestellt, deren Wurzeln $\psi_1^2$, $\psi_2^2$, $\psi_3^2$ sind, und haben auch ihre Discriminante gebildet. Im vorigen Paragraphen haben wir ein Kennzeichen für die Anzahl der positiven Wurzeln einer cubischen Gleichung kennen gelernt, woraus sich folgendes Resultat ergiebt:

Ist $D < 0$, so hat die biquadratische Gleichung zwei reelle und zwei imaginäre Wurzeln.

Ist $D > 0$, so müssen, wenn vier reelle Wurzeln vorhanden sein sollen, in der Reihe
\[
1, \quad H, \quad H^2 - 16 A f^2, \quad -T^2
\]
drei Zeichenwechsel vorkommen, d. h. es muss
\[
H < 0, \quad H^2 - 16 A f^2 > 0
\]
sein. Bei den drei anderen noch möglichen Zeichencombinationen sind alle vier Wurzeln imaginär. Hierbei aber kann für $x$ ein beliebiger reeller Werth gesetzt werden.
